% GENERATED FILE. Edit canonical lessons, then run npm run generate:books.
% canonical-source-hash: fnv1a64:4939a50a971712d2
% canonical-lessons: PT-C529-detergente, PT-C529-estender, PT-C529-estendal, PT-C529-mola, PT-C529-engomar

\chapter{A roupa lavada}
\label{ch:pt-a-roupa-lavada}
\hlchaptermodality{529}

\begin{chapteropening}
\textbf{By the end of this chapter:} \emph{I can talk about hanging out the washing and ironing.}

Complete the last lesson of chapter 529: I can talk about hanging out the washing and ironing.
\end{chapteropening}

\section[o detergente]{o detergente — detergent, washing powder}

\label{lesson:PT-C529-detergente}



\begin{quote}
Before the new one: say the Portuguese for to scrub, then the Portuguese for mop.
\end{quote}

\subsection*{You'll want to know: o detergente}
\textbf{o detergente} — \textquotedblleft{}detergent, washing powder\textquotedblright{}.

\begin{quote}
Acabou o detergente da máquina.
\end{quote}

\begin{quote}
\emph{We have run out of washing powder.}
\end{quote}

\emph{Detergente da loiça} — washing-up liquid.

\begin{cousinweb}
\emph{detergente} is a learned word built on Latin \emph{dētergēns} \textquotedblleft{}wiping off\textquotedblright{}, from \emph{dētergēre}, from \emph{tergēre} \textquotedblleft{}to wipe\textquotedblright{} (Lewis \& Short).
\end{cousinweb}

\begin{grammarlens}[title={What you've built}]
One more word for telling, explaining and describing at B1.
\end{grammarlens}

\subsection*{Your turn}
\begin{itemize}
\raggedright
  \item \emph{Say it:} \emph{o detergente}
  \item \emph{Say it:} \emph{o detergente}, once more
  \item \emph{Say it:} a sentence of your own with \emph{o detergente}
  \item \emph{Recall:} say the Portuguese for to scrub, then the Portuguese for mop, then say \emph{o detergente} again
\end{itemize}

\begin{tcolorbox}[breakable,colback=teal!4,colframe=teal!35!black,title={Before you move on}]
What does o detergente mean? (\textquotedblleft{}Detergent, washing powder\textquotedblright{}.) Say it once more.
\end{tcolorbox}
\section[estender]{estender — to hang out (washing); to stretch out}

\label{lesson:PT-C529-estender}



\begin{quote}
Before the new one: say the Portuguese for mop, then the Portuguese for detergent.
\end{quote}

\subsection*{You'll want to know: estender}
\textbf{estender} — \textquotedblleft{}to hang out (washing); to stretch out\textquotedblright{}.

\begin{quote}
Estendo a roupa ao sol.
\end{quote}

\begin{quote}
\emph{I hang the washing out in the sun.}
\end{quote}

\emph{Estender a roupa} — to hang out the washing.

\begin{cousinweb}
\emph{estender} is inherited, through Old Portuguese, from Latin \emph{extendere} \textquotedblleft{}to stretch out\textquotedblright{} (EtymDB), built on \emph{tendere} \textquotedblleft{}to stretch\textquotedblright{}.
\end{cousinweb}

\begin{grammarlens}[title={What you've built}]
One more word for telling, explaining and describing at B1.
\end{grammarlens}

\subsection*{Your turn}
\begin{itemize}
\raggedright
  \item \emph{Say it:} \emph{estender}
  \item \emph{Say it:} \emph{estender}, once more
  \item \emph{Say it:} a sentence of your own with \emph{estender}
  \item \emph{Recall:} say the Portuguese for mop, then the Portuguese for detergent, then say \emph{estender} again
\end{itemize}

\begin{tcolorbox}[breakable,colback=teal!4,colframe=teal!35!black,title={Before you move on}]
What does estender mean? (\textquotedblleft{}To hang out (washing); to stretch out\textquotedblright{}.) Say it once more.
\end{tcolorbox}
\section[o estendal]{o estendal — clothes horse; washing line}

\label{lesson:PT-C529-estendal}



\begin{quote}
Before the new one: say the Portuguese for detergent, then the Portuguese for to hang out.
\end{quote}

\subsection*{You'll want to know: o estendal}
\textbf{o estendal} — \textquotedblleft{}clothes horse; washing line\textquotedblright{}.

\begin{quote}
O estendal está na varanda.
\end{quote}

\begin{quote}
\emph{The clothes horse is on the balcony.}
\end{quote}

European Portuguese; in Brazil, \emph{o varal}.

\begin{cousinweb}
\emph{estendal} is built on \emph{estender} \textquotedblleft{}to hang out, to stretch\textquotedblright{}, inherited from Latin \emph{extendere} (EtymDB).
\end{cousinweb}

\begin{grammarlens}[title={What you've built}]
One more word for telling, explaining and describing at B1.
\end{grammarlens}

\subsection*{Your turn}
\begin{itemize}
\raggedright
  \item \emph{Say it:} \emph{o estendal}
  \item \emph{Say it:} \emph{o estendal}, once more
  \item \emph{Say it:} a sentence of your own with \emph{o estendal}
  \item \emph{Recall:} say the Portuguese for detergent, then the Portuguese for to hang out, then say \emph{o estendal} again
\end{itemize}

\begin{tcolorbox}[breakable,colback=teal!4,colframe=teal!35!black,title={Before you move on}]
What does o estendal mean? (\textquotedblleft{}Clothes horse; washing line\textquotedblright{}.) Say it once more.
\end{tcolorbox}
\section[a mola]{a mola — clothes peg; spring}

\label{lesson:PT-C529-mola}



\begin{quote}
Before the new one: say the Portuguese for to hang out, then the Portuguese for clothes horse.
\end{quote}

\subsection*{You'll want to know: a mola}
\textbf{a mola} — \textquotedblleft{}clothes peg; spring\textquotedblright{}.

\begin{quote}
Prendo as meias com molas.
\end{quote}

\begin{quote}
\emph{I peg the socks up with clothes pegs.}
\end{quote}

Also a spring: \emph{a mola do colchão}.

\begin{cousinweb}
\emph{mola} was borrowed from Italian \emph{molla} \textquotedblleft{}a spring\textquotedblright{} (EtymDB), which goes back to Latin \emph{mollis} \textquotedblleft{}soft, pliant\textquotedblright{}.
\end{cousinweb}

\begin{grammarlens}[title={What you've built}]
One more word for telling, explaining and describing at B1.
\end{grammarlens}

\subsection*{Your turn}
\begin{itemize}
\raggedright
  \item \emph{Say it:} \emph{a mola}
  \item \emph{Say it:} \emph{a mola}, once more
  \item \emph{Say it:} a sentence of your own with \emph{a mola}
  \item \emph{Recall:} say the Portuguese for to hang out, then the Portuguese for clothes horse, then say \emph{a mola} again
\end{itemize}

\begin{tcolorbox}[breakable,colback=teal!4,colframe=teal!35!black,title={Before you move on}]
What does a mola mean? (\textquotedblleft{}Clothes peg; spring\textquotedblright{}.) Say it once more.
\end{tcolorbox}
\section[engomar]{engomar — to iron; to starch}

\label{lesson:PT-C529-engomar}



\begin{quote}
Before the new one: say the Portuguese for clothes horse, then the Portuguese for clothes peg.
\end{quote}

\subsection*{You'll want to know: engomar}
\textbf{engomar} — \textquotedblleft{}to iron; to starch\textquotedblright{}.

\begin{quote}
O meu pai engoma as camisas ao domingo.
\end{quote}

\begin{quote}
\emph{My father irons the shirts on Sundays.}
\end{quote}

Also \emph{passar a ferro}.

\begin{cousinweb}
\emph{engomar} is built on \emph{goma} \textquotedblleft{}starch, gum\textquotedblright{}, which goes back to Latin \emph{gummi}, from Greek \emph{kómmi}; the Greek word is usually traced to Egyptian.
\end{cousinweb}

\begin{grammarlens}[title={What you've built}]
One more word for telling, explaining and describing at B1.
\end{grammarlens}

\subsection*{Your turn}
\begin{itemize}
\raggedright
  \item \emph{Say it:} \emph{engomar}
  \item \emph{Say it:} \emph{engomar}, once more
  \item \emph{Say it:} a sentence of your own with \emph{engomar}
  \item \emph{Recall:} say the Portuguese for clothes horse, then the Portuguese for clothes peg, then say \emph{engomar} again
\end{itemize}

\begin{tcolorbox}[breakable,colback=teal!4,colframe=teal!35!black,title={Before you move on}]
What does engomar mean? (\textquotedblleft{}To iron; to starch\textquotedblright{}.) Say it once more.
\end{tcolorbox}
